\chapter{Bariatric Surgery}

\section*{What Is This Surgery?}
Bariatric surgery, also known as metabolic or weight-loss surgery, encompasses a range of surgical procedures that modify the gastrointestinal tract to induce significant and durable weight loss, leading to the improvement or resolution of obesity-related medical conditions. These procedures primarily target the stomach and small intestine, altering their anatomy and physiology to reduce food intake, decrease nutrient absorption, and profoundly impact gut hormone signaling. The goal is not merely cosmetic weight reduction but rather the comprehensive treatment of severe obesity as a chronic disease, addressing its complex metabolic consequences such as type 2 diabetes, hypertension, and dyslipidemia. By achieving substantial and sustained weight loss, bariatric surgery significantly improves overall health, quality of life, and life expectancy for eligible patients.

\section*{Alternative Names}
\begin{itemize}
  \item Weight-loss surgery
  \item Metabolic surgery
  \item Gastric bypass (specifically Roux-en-Y Gastric Bypass, RYGB)
  \item Sleeve gastrectomy (specifically Laparoscopic Sleeve Gastrectomy, LSG)
  \item Biliopancreatic Diversion with Duodenal Switch (BPD/DS)
\end{itemize}

\section*{Relevant Surgical Anatomy}
Bariatric surgery fundamentally alters the anatomy of the upper gastrointestinal tract, primarily involving the stomach and small intestine. The stomach is a J-shaped organ with distinct regions: the cardia, fundus, body, antrum, and pylorus. Its lesser curvature is supplied by the left and right gastric arteries, while the greater curvature receives blood from the left and right gastroepiploic arteries, all branches of the celiac axis. The vagus nerve provides parasympathetic innervation, influencing gastric motility and acid secretion. The small intestine, extending from the pylorus to the ileocecal valve, is divided into the duodenum, jejunum, and ileum. The ligament of Treitz marks the duodenojejunal flexure and is a crucial landmark for identifying the jejunum.

For a sleeve gastrectomy, the surgical focus is on the greater curvature of the stomach, requiring meticulous division of the short gastric vessels within the gastrosplenic ligament to free the fundus and body. The staple line typically runs parallel to the lesser curvature, preserving the antrum and pylorus. In Roux-en-Y gastric bypass, a small gastric pouch is created from the proximal stomach, necessitating division of the stomach near the angle of His. The jejunum is then transected distal to the ligament of Treitz to create the Roux limb (alimentary limb) and the biliopancreatic limb. The mesenteric defects created by rerouting the bowel must be carefully closed to prevent internal hernias. Understanding the precise vascular supply and nerve distribution is critical to minimize complications such as bleeding, ischemia, or vagal nerve injury.

\section{Surgical Principle \& Rationale}
The fundamental principle of bariatric surgery is to induce a negative energy balance and favorable metabolic changes through a combination of restrictive, malabsorptive, and neurohormonal mechanisms. Restrictive procedures, such as sleeve gastrectomy, reduce the functional volume of the stomach, limiting the amount of food that can be consumed at one time and promoting early satiety. This physical restriction is complemented by the removal of the gastric fundus, which is the primary site of ghrelin production, leading to a reduction in appetite-stimulating hormones.

Malabsorptive components, primarily seen in Roux-en-Y gastric bypass and biliopancreatic diversion with duodenal switch, involve rerouting the small intestine to bypass significant lengths of the digestive tract, thereby reducing the absorption of calories and nutrients. Beyond these mechanical effects, a crucial aspect of bariatric surgery's efficacy lies in its profound impact on gut hormone physiology. Procedures like gastric bypass and sleeve gastrectomy alter the secretion patterns of various enteroendocrine hormones, including glucagon-like peptide-1 (GLP-1), peptide YY (PYY), and cholecystokinin (CCK). These hormonal shifts contribute to improved glucose homeostasis, reduced appetite, increased satiety, and altered energy expenditure, leading to the resolution of metabolic diseases like type 2 diabetes independently of weight loss. The technical goals of these operations include creating a small, functional gastric pouch, ensuring secure anastomoses or staple lines, and minimizing long-term complications while maximizing metabolic benefits.

\section{History \& Surgical Evolution}
The origins of bariatric surgery can be traced back to the 1950s with the development of the jejunoileal bypass by Kremen and Linner, a purely malabsorptive procedure that, while effective for weight loss, was fraught with severe metabolic complications such as liver failure and electrolyte disturbances, leading to its eventual abandonment. In the 1960s, Dr. Edward E. Mason pioneered the gastric bypass, initially as a horizontal gastroplasty, marking a shift towards procedures that combined restriction with some degree of malabsorption. The modern Roux-en-Y gastric bypass, involving a small gastric pouch and a Roux limb, was refined over the subsequent decades, becoming the gold standard for many years.

The 1990s witnessed a revolutionary shift with the advent of laparoscopic surgery. Dr. Alan Wittgrove performed the first laparoscopic Roux-en-Y gastric bypass in 1993, transforming bariatric surgery from a highly invasive open procedure to a minimally invasive approach, significantly reducing patient morbidity, pain, and recovery time. The early 2000s saw the rise in popularity of the laparoscopic sleeve gastrectomy, initially as a component of the biliopancreatic diversion with duodenal switch, but later recognized as a standalone procedure by Dr. Michel Gagner due to its simplicity, efficacy, and favorable risk profile. Adjustable gastric banding also gained traction in the late 1990s and early 2000s but has since seen a decline in use due to lower efficacy and higher long-term complication rates compared to sleeve gastrectomy and gastric bypass.

\section{Clinical Indications}
Bariatric surgery is indicated for individuals who meet specific criteria for severe obesity and have failed to achieve sustainable weight loss through non-surgical methods. The primary indications, as outlined by national and international guidelines, include:

\begin{itemize}
  \item A Body Mass Index (BMI) of 40 kg/m\textasciicircum2 or greater (Class III obesity), regardless of the presence of obesity-related comorbidities.
  \item A BMI of 35 kg/m\textasciicircum2 to 39.9 kg/m\textasciicircum2 (Class II obesity) with at least one significant obesity-related comorbidity, such as type 2 diabetes, hypertension, or severe obstructive sleep apnea.
  \item A BMI of 30 kg/m\textasciicircum2 to 34.9 kg/m\textasciicircum2 (Class I obesity) with uncontrolled type 2 diabetes or other severe metabolic diseases, particularly when medical management is insufficient.
  \item Type 2 diabetes mellitus that is difficult to control with conventional medical therapy, often leading to significant end-organ damage despite optimal pharmacological intervention.
  \item Severe hypertension requiring multiple medications for control or associated with end-organ damage, such as left ventricular hypertrophy or nephropathy.
  \item Obstructive sleep apnea (OSA) requiring continuous positive airway pressure (CPAP) or associated with significant daytime somnolence and cardiovascular risk.
  \item Non-alcoholic fatty liver disease (NAFLD) or non-alcoholic steatohepatitis (NASH) with evidence of fibrosis, which can progress to cirrhosis.
  \item Significant osteoarthritis of weight-bearing joints (e.g., knees, hips), causing debilitating pain and limiting mobility, where weight loss can significantly improve symptoms and delay joint replacement.
\end{itemize}

\section{Clinical Positioning}
For individuals with severe obesity and its associated metabolic complications, bariatric surgery is increasingly recognized as a primary, definitive therapy due to its superior efficacy in achieving sustained weight loss and disease remission compared to non-surgical interventions. While lifestyle modifications and pharmacotherapy are often attempted first, their long-term success rates for severe obesity are generally low, with high rates of weight regain. Therefore, for eligible patients, bariatric surgery is not merely a last resort but a highly effective first-line treatment option that can significantly alter the disease trajectory, reduce medication burden, and improve life expectancy.

In some cases, particularly for patients with a BMI of 30-34.9 kg/m\textasciicircum2 and uncontrolled type 2 diabetes, it may even be considered a primary intervention alongside intensive medical therapy, as supported by trials like STAMPEDE. Bariatric surgery can also serve as a secondary or revisional procedure for patients who have experienced inadequate weight loss or significant complications from a previous bariatric operation, such as conversion of an adjustable gastric band to a sleeve gastrectomy or Roux-en-Y gastric bypass, or revision of a failed sleeve gastrectomy to a bypass.

\section{Surgical Technique \& Procedure}
Bariatric surgery is almost exclusively performed laparoscopically, utilizing several small incisions (typically 0.5-1.5 cm) rather than a large open incision, which minimizes pain, reduces recovery time, and lowers the risk of wound complications. The procedure is carried out under general anesthesia. The patient is typically positioned supine, and the abdomen is insufflated with carbon dioxide to create a pneumoperitoneum and a working space. Trocars are then inserted to allow the passage of a camera and specialized surgical instruments.

\begin{itemize}
  \item \textbf{Laparoscopic Sleeve Gastrectomy (LSG):}
    \begin{enumerate}
      \item \textbf{Anesthesia \& Positioning:} General anesthesia is induced. The patient is positioned supine with arms tucked, often in a reverse Trendelenburg position to facilitate visualization of the upper abdomen.
      \item \textbf{Access \& Insufflation:} A Veress needle or optical trocar is used to gain abdominal access, followed by insufflation with CO\textsubscript{2} to 15 mmHg. Multiple trocars (typically 4-5) are placed in the upper abdomen.
      \item \textbf{Liver Retraction:} The left lobe of the liver is retracted superiorly to expose the gastroesophageal junction and fundus.
      \item \textbf{Gastric Mobilization:} The greater curvature of the stomach is mobilized by dividing the gastrosplenic ligament and short gastric vessels, starting 4-6 cm from the pylorus and extending up to the angle of His, freeing the stomach from the spleen and surrounding attachments.
      \item \textbf{Bougie Placement:} A calibration bougie (typically 32-40 French) is inserted orally into the stomach and guided along the lesser curvature to serve as a template for the new stomach size.
      \item \textbf{Gastric Resection:} A surgical stapler is then used to divide the stomach longitudinally, starting approximately 4-6 cm from the pylorus and extending up towards the angle of His, removing about 75-80\% of the stomach. The resected portion is extracted.
      \item \textbf{Staple Line Reinforcement \& Leak Test:} The staple line is often reinforced with sutures or buttressing material to reduce the risk of bleeding and leak. A leak test is performed by inflating the new stomach with air or methylene blue solution while submerged in saline.
      \item \textbf{Closure:} Trocars are removed under direct visualization, and fascia at larger port sites is closed. Skin incisions are closed with absorbable sutures or staples. Drains are rarely used routinely.
    \end{enumerate}

  \item \textbf{Laparoscopic Roux-en-Y Gastric Bypass (RYGB):}
    \begin{enumerate}
      \item \textbf{Anesthesia \& Positioning:} Similar to LSG, general anesthesia and supine/reverse Trendelenburg positioning.
      \item \textbf{Access \& Insufflation:} Similar trocar placement, often with an additional port for liver retraction.
      \item \textbf{Gastric Pouch Creation:} A small gastric pouch (approximately 15-30 mL) is created by dividing the stomach horizontally and vertically, separating it from the larger remnant stomach using surgical staplers.
      \item \textbf{Jejunal Transection:} The jejunum is identified and transected approximately 50-100 cm distal to the ligament of Treitz.
      \item \textbf{Roux Limb Creation \& Gastrojejunostomy:} The distal segment of the jejunum (Roux limb) is brought up in an antecolic or retrocolic fashion and connected to the newly created gastric pouch, forming the gastrojejunostomy (typically hand-sewn or stapled).
      \item \textbf{Jejunojejunostomy:} The proximal segment of the jejunum (biliopancreatic limb), which carries bile and pancreatic enzymes, is then connected to the Roux limb approximately 75-150 cm distal to the gastrojejunostomy, forming the jejunojejunostomy.
      \item \textbf{Mesenteric Defect Closure:} The mesenteric defects created by rerouting the bowel (Petersen's defect and the jejunojejunostomy mesenteric defect) are meticulously closed with non-absorbable sutures to prevent internal hernias.
      \item \textbf{Leak Test \& Closure:} A leak test is performed at both anastomoses. Trocars are removed, fascia closed, and skin incisions are closed. Drains are typically not placed routinely.
    \end{enumerate}
\end{itemize}

\section{Preoperative Workup \& Preparation}
Thorough preoperative preparation is crucial for optimizing outcomes and minimizing risks in bariatric surgery. This typically involves a multidisciplinary evaluation over several weeks or months to ensure the patient is medically and psychologically ready for surgery and the lifelong commitment required.

\begin{itemize}
  \item \textbf{Medical Evaluation:} Comprehensive assessment by the bariatric surgeon, internist, and often an endocrinologist. This includes a detailed medical history, physical examination, and extensive laboratory testing (complete blood count, metabolic panel, liver function tests, thyroid function tests, lipid panel, hemoglobin A1c, vitamin D, B12, folate, iron studies).
  \item \textbf{Cardiac and Pulmonary Clearance:} Electrocardiogram (ECG), and often a cardiac stress test or echocardiogram, especially for patients with known cardiac disease or significant risk factors. Pulmonary function tests and a sleep study may be required for patients with respiratory symptoms or suspected obstructive sleep apnea.
  \item \textbf{Gastrointestinal Evaluation:} Upper endoscopy (EGD) is routinely performed to assess for hiatal hernia, Helicobacter pylori infection (which is treated if positive), Barrett's esophagus, or other gastric pathologies. Abdominal ultrasound may be performed to evaluate for gallstones, which may prompt prophylactic cholecystectomy in some centers.
  \item \textbf{Nutritional Counseling:} Extensive education by a registered dietitian on dietary changes, vitamin and mineral supplementation, and long-term eating habits. Patients often undergo a very low-calorie diet (VLCD) for 2-4 weeks preoperatively to reduce liver size, making the surgery safer and technically easier.
  \item \textbf{Psychological Evaluation:} Assessment by a psychologist or psychiatrist to identify and address any untreated mental health conditions (e.g., severe depression, active substance abuse, uncontrolled eating disorders) that could impact surgical outcomes or compliance with postoperative care.
  \item \textbf{Lifestyle Modifications:} Patients are strongly advised to stop smoking and abstain from alcohol for a specified period before surgery. Regular physical activity is encouraged.
  \item \textbf{Medication Adjustments:} Review and adjustment of medications, particularly anticoagulants, antiplatelet agents, and diabetes medications, to ensure safe perioperative management.
  \item \textbf{Informed Consent:} Detailed discussion of the procedure, potential benefits, risks, and lifelong commitment required for success.
  \item \textbf{Preoperative Antibiotics:} Administered intravenously shortly before incision to reduce the risk of surgical site infection.
  \item \textbf{NPO Status:} Patients are instructed to be NPO (nil per os) after midnight on the day of surgery.
\end{itemize}

\section{Contraindications \& Precautions}
Careful patient selection is paramount to ensure safety and optimize outcomes in bariatric surgery.

\textbf{Absolute Contraindications:}
\begin{itemize}
  \item Active malignancy or a limited life expectancy due to other terminal illnesses (e.g., end-stage organ failure).
  \item Severe, untreated psychiatric illness (e.g., uncontrolled psychosis, severe depression with suicidal ideation) that would impair the patient's ability to understand or comply with postoperative care.
  \item Active substance abuse (alcohol or illicit drugs) within a specified period (e.g., 6-12 months) prior to surgery, as it compromises adherence and increases complication risk.
  \item Inability or unwillingness to comply with lifelong dietary restrictions, vitamin supplementation, and regular follow-up appointments.
  \item Pregnancy or planning pregnancy within 12-18 months post-surgery due to rapid weight loss and potential nutritional deficiencies.
  \item Uncontrolled coagulopathy or bleeding disorder that cannot be safely managed perioperatively.
\end{itemize}

\textbf{Relative Contraindications and Precautions:}
\begin{itemize}
  \item Severe cardiac or pulmonary disease that significantly increases anesthetic and surgical risk, requiring extensive preoperative optimization and careful risk-benefit assessment.
  \item Portal hypertension with esophageal varices (especially for sleeve gastrectomy due to potential for staple line bleeding and exacerbation of varices).
  \item Inflammatory bowel disease (e.g., Crohn's disease) may complicate malabsorptive procedures and increase the risk of anastomotic complications.
  \item Advanced age with significant frailty, requiring individualized risk-benefit assessment and consideration of less invasive options.
  \item Certain eating disorders (e.g., bulimia nervosa) that may be exacerbated by surgical changes, necessitating intensive psychological support.
  \item Chronic use of non-steroidal anti-inflammatory drugs (NSAIDs) due to increased risk of marginal ulcers after gastric bypass; alternative pain management strategies are crucial.
  \item Significant hiatal hernia, which may require repair during sleeve gastrectomy or favor gastric bypass in patients with severe, refractory gastroesophageal reflux disease (GERD).
\end{itemize}

\section{Complications \& Risks}
While generally safe, bariatric surgery carries potential risks, both immediate and long-term, which are carefully discussed with patients during the informed consent process.

\textbf{General Surgical and Anesthesia Risks:}
\begin{itemize}
  \item \textbf{Anesthesia complications:} Adverse reactions to medications, respiratory problems (e.g., hypoventilation, atelectasis), cardiac events (e.g., myocardial infarction, stroke), aspiration pneumonia.
  \item \textbf{Bleeding:} Intraoperative or postoperative hemorrhage, which may manifest as hematemesis, melena, or intra-abdominal collection, potentially requiring blood transfusion or reoperation.
  \item \textbf{Infection:} Surgical site infection (superficial wound infection, deep space infection, intra-abdominal abscess), pneumonia, urinary tract infection, or sepsis.
  \item \textbf{Thromboembolic events:} Deep vein thrombosis (DVT) in the lower extremities and pulmonary embolism (PE), which can be life-threatening. Prophylaxis with anticoagulants and early ambulation is routinely used.
  \item \textbf{Injury to adjacent organs:} Accidental damage to the spleen, liver, bowel, or other abdominal structures during trocar insertion or dissection.
\end{itemize}

\textbf{Procedure-Specific Risks (Immediate and Long-term):}
\begin{itemize}
  \item \textbf{Anastomotic leak (RYGB) or staple line leak (LSG):} A serious and potentially life-threatening complication where digestive contents leak from the surgical connections or staple line, leading to peritonitis, sepsis, and potentially multiple reoperations, prolonged hospitalization, or even death.
  \item \textbf{Anastomotic stricture (RYGB) or staple line stricture (LSG):} Narrowing of the surgical connection or stomach outlet, causing difficulty swallowing (dysphagia), nausea, and vomiting, often requiring endoscopic dilation.
  \item \textbf{Marginal ulcer (RYGB):} Ulcers that can form at the gastrojejunostomy site, causing pain, bleeding, or perforation, often associated with NSAID use, smoking, or H. pylori.
  \item \textbf{Dumping syndrome (RYGB):} Rapid emptying of undigested food, particularly high-sugar or high-fat items, into the small intestine, leading to symptoms like nausea, vomiting, diarrhea, abdominal cramps, sweating, and dizziness.
  \item \textbf{Internal hernia (RYGB):} A serious complication where a loop of bowel herniates through a defect in the mesentery, leading to bowel obstruction, ischemia, and potentially necrosis, requiring urgent surgical intervention.
  \item \textbf{Nutritional deficiencies:} Long-term deficiencies in vitamins (e.g., B12, D, folate, thiamine), minerals (e.g., iron, calcium, zinc, copper), and protein, requiring lifelong supplementation and monitoring.
  \item \textbf{Gallstones:} Rapid weight loss increases the risk of gallstone formation, sometimes requiring cholecystectomy.
  \item \textbf{Gastroesophageal reflux disease (GERD) exacerbation (LSG):} While some patients with GERD improve after LSG, others may develop new or worsened reflux symptoms, occasionally requiring conversion to RYGB.
  \item \textbf{Weight regain:} Despite initial success, some patients may experience significant weight regain over time, often due to non-adherence to dietary and lifestyle recommendations, or maladaptive eating behaviors.
\end{itemize}

\section{Postoperative Recovery Timeline}
Immediately after bariatric surgery, patients are transferred to the Post-Anesthesia Care Unit (PACU) for close monitoring of vital signs, pain control, and recovery from anesthesia. Early ambulation is encouraged within hours of surgery to prevent deep vein thrombosis and pulmonary complications. Most patients remain hospitalized for 1 to 3 days, depending on the specific procedure and their individual recovery progress. During this initial period, pain is managed with intravenous or oral analgesics, and patients begin a clear liquid diet, gradually progressing to full liquids. Nurses and dietitians provide extensive education on the new dietary regimen, emphasizing small, frequent meals and adequate hydration.

Over the first 1-2 weeks post-discharge, patients adhere strictly to a liquid diet, followed by a pureed diet for several weeks. Nutritional supplements, including protein shakes and chewable multivitamins, are initiated. Patients are advised to avoid heavy lifting (typically over 10-15 lbs) and strenuous activity for 4-6 weeks to allow incision healing, but are strongly encouraged to engage in light walking and gradually increase activity. Wound care instructions are provided, and any drains, if placed, are typically removed before discharge or at the first postoperative visit. Long-term recovery involves a gradual progression to soft foods and then regular, healthy solid foods, with a strong emphasis on lean protein, non-starchy vegetables, and avoidance of high-sugar and high-fat items. Lifelong follow-up with the bariatric team, including regular clinic visits, laboratory monitoring, and adherence to vitamin and mineral supplementation, is essential for sustained weight loss, prevention of nutritional deficiencies, and management of any potential complications. Weight loss is typically most rapid in the first 6-12 months, with continued loss for up to 18-24 months, followed by a plateau and maintenance phase.

\section{Surgical Findings \& Pathology}
During bariatric surgery, the surgeon documents key intraoperative findings, which may include the size of the liver (often enlarged due to steatosis, necessitating preoperative diet), presence and size of any hiatal hernia, adhesions from previous surgeries, and the overall condition of the abdominal organs. For sleeve gastrectomy, the resected stomach specimen is sent for pathological examination. The pathologist confirms the presence of gastric tissue, assesses for any incidental findings such as Helicobacter pylori infection, gastritis, polyps, or rarely, early neoplastic changes. For gastric bypass, any suspicious lesions identified during preoperative endoscopy or intraoperatively may be biopsied and sent for pathology. The pathology report provides a definitive histological diagnosis of the resected or biopsied tissues, which is crucial for guiding postoperative management and surveillance. While the primary goal of bariatric surgery is metabolic, the pathological examination ensures no underlying gastric pathology is missed and confirms the anatomical changes made.

\section{Expected Postoperative Course}
A normal or successful postoperative course following bariatric surgery is characterized by several key indicators. These include uncomplicated wound healing without signs of infection or dehiscence, the absence of anastomotic or staple line leaks, and no evidence of strictures, internal hernias, or other acute surgical complications. Patients typically experience progressive and sustained weight loss, achieving significant reduction in excess body weight within the first 1-2 years. Crucially, there is a resolution or marked improvement in obesity-related comorbidities such as type 2 diabetes, hypertension, and obstructive sleep apnea, often leading to reduced medication requirements or complete discontinuation. Patients should tolerate their prescribed diet progression well, without persistent nausea, vomiting, severe dumping syndrome, or intractable pain. Furthermore, laboratory monitoring should show stable micronutrient levels, indicating effective adherence to supplementation and preventing deficiencies. Ultimately, a normal result signifies an improved quality of life, enhanced physical function, and a reduced risk of obesity-related morbidity and mortality, allowing patients to engage more fully in daily activities and enjoy a healthier lifestyle.

\section{Postoperative Care \& Follow-up}
Lifelong follow-up is a cornerstone of successful bariatric surgery outcomes, ensuring sustained weight loss, prevention of nutritional deficiencies, and early detection and management of complications.

\begin{itemize}
  \item \textbf{Clinic Visits:} Regular postoperative visits are scheduled frequently in the first year (e.g., 1 week, 1 month, 3 months, 6 months, 12 months) and annually thereafter with the bariatric surgeon, dietitian, and often a bariatric nurse or physician assistant.
  \item \textbf{Laboratory Monitoring:} Lifelong annual blood tests are essential to screen for nutritional deficiencies. These typically include a complete blood count (CBC), comprehensive metabolic panel, liver function tests, hemoglobin A1c, lipid panel, and specific vitamin and mineral levels (e.g., vitamin B12, vitamin D, folate, iron studies, calcium, parathyroid hormone, zinc, copper, thiamine).
  \item \textbf{Bone Density Evaluation:} Dual-energy X-ray absorptiometry (DEXA) scans are recommended periodically (e.g., every 1-2 years) to monitor for bone loss, particularly in patients at higher risk for osteoporosis due to rapid weight loss and potential calcium/vitamin D malabsorption.
  \item \textbf{Imaging Studies:} Abdominal ultrasound may be performed if gallstone symptoms develop. Upper endoscopy may be indicated for symptoms such as persistent nausea, vomiting, dysphagia, or suspicion of marginal ulcer or stricture.
  \item \textbf{Physical Activity and Rehabilitation:} Patients are encouraged to gradually increase physical activity, often with guidance from a physical therapist or exercise physiologist, to improve strength, mobility, and overall fitness.
  \item \textbf{Psychological Support:} Ongoing access to psychological counseling or support groups can be beneficial for managing body image changes, emotional eating, and adapting to new lifestyle demands.
  \item \textbf{Warning Signs:} Patients are educated to seek immediate medical attention for warning signs such as persistent fever (above 101.5\textdegree F), severe or worsening abdominal pain, persistent nausea or vomiting, inability to tolerate liquids, rapid heart rate, shortness of breath, or changes in wound appearance (redness, swelling, purulent discharge).
\end{itemize}
Adherence to these comprehensive follow-up recommendations is critical for long-term health, prevention of complications, and maximizing the benefits of bariatric surgery.

\section{Epidemiology}
Obesity is a global epidemic, with its prevalence steadily increasing worldwide. In the United States, over 40\% of adults are classified as obese, and approximately 9\% have severe obesity (BMI \textgreater{}40 kg/m\textasciicircum2). This rising prevalence has led to a corresponding increase in the demand for bariatric surgery as an effective treatment modality. Annually, over 250,000 bariatric procedures are performed in the U.S., with laparoscopic sleeve gastrectomy and Roux-en-Y gastric bypass being the most common types, accounting for over 90\% of all operations. The typical bariatric surgery patient is often female, in their mid-40s, and presents with multiple obesity-related comorbidities, with type 2 diabetes, hypertension, and obstructive sleep apnea being the most frequent indications. While the overall mortality rate for bariatric surgery is remarkably low, typically less than 0.3\% in experienced centers, the morbidity rate, encompassing complications such as leaks, strictures, and infections, ranges from 5-10\%. Despite these risks, bariatric surgery significantly reduces overall mortality from obesity-related causes and can increase life expectancy by several years, making it a highly impactful public health intervention.

\section{Outcomes \& Prognosis}
The outcomes and prognosis following bariatric surgery are generally excellent, offering significant and sustained improvements in health and quality of life for the vast majority of patients. Typical success rates for achieving at least 50\% excess weight loss range from 60-80\% depending on the procedure type and patient adherence, with peak weight loss usually occurring within 18-24 months. Beyond weight loss, the most profound impact is on metabolic health: type 2 diabetes remission rates are 60-80\% for RYGB and 40-60\% for LSG, often occurring rapidly post-surgery, as demonstrated by studies like the STAMPEDE trial. Hypertension and dyslipidemia also show significant improvement or resolution in 70-80\% of patients. Functional outcomes include improved mobility, reduced joint pain, resolution of obstructive sleep apnea, and enhanced overall physical and mental well-being. Long-term studies, such as the Swedish Obese Subjects (SOS) study, demonstrate a significant reduction in cardiovascular events, cancer risk, and overall mortality compared to severely obese individuals who do not undergo surgery. While weight regain is a possibility (around 15-20\% of initial weight loss over 5-10 years), it is typically less severe than with non-surgical interventions. Prognostic factors for optimal outcomes include younger age, lower preoperative BMI, absence of severe comorbidities, and strict adherence to postoperative dietary guidelines, vitamin supplementation, and regular follow-up.

\section{References}
\begin{enumerate}
  \item MedlinePlus. Bariatric surgery. U.S. National Library of Medicine; 2025.
  \item American Society for Metabolic and Bariatric Surgery (ASMBS). Bariatric Surgery Procedures. Available at: www.asmbs.org; 2024.
  \item Sabiston Textbook of Surgery: The Biological Basis of Modern Surgical Practice. 22nd ed. Elsevier; 2022.
  \item Schauer PR, Bhatt DL, Kirwan JF, et al. Bariatric Surgery versus Intensive Medical Therapy for Diabetes -- 5-Year Outcomes. N Engl J Med. 2017;376(7):641-651.
\end{enumerate}

\clearpage